\section{طرح الإشكالية}
تُعدّ الأسرة أوّل بيئة اجتماعية ينشأ فيها الفرد، ومنها يكتسب أنماط السلوك والقيم وأساليب التفاعل
مع ذاته ومع الآخرين. ويرى حامد عبد السلام زهران أنّ الأسرة تمثّل الجماعة الاجتماعية الأولى المسؤولة
عن عملية التنشئة الاجتماعية، إذ تؤدّي دورها من خلال أساليب الضبط والثواب والعقاب، بما يسهم في
تشكيل شخصية الفرد وبناء نظرته إلى نفسه وإلى الآخرين (زهران، 1986). ولا يقتصر أثر هذه
الممارسات على مرحلة الطفولة، بل يمتدّ إلى مرحلة المراهقة التي تُعدّ من أكثر المراحل حساسيةً في النموّ
الإنساني، لما تشهده من تغيّرات جسمية ومعرفية وانفعالية واجتماعية متسارعة، يصاحبها سعي المراهق إلى
تكوين هويّته وتحقيق استقلاليته وبناء صورة أكثر وضوحًا عن ذاته.

وفي هذا السياق، تكتسب أساليب المعاملة الوالدية أهميةً خاصّة، إذ تمثّل أحد أهمّ العوامل الأسرية التي
تسهم في تشكيل خبرات المراهق النفسية والاجتماعية. غير أنّ الأثر النفسي لهذه الأساليب لا يتحدّد
بالممارسات الوالدية في حدّ ذاتها، وإنما يتأثّر كذلك بالطريقة التي يدرك بها المراهق تلك الممارسات
ويفسّرها في ضوء خبراته وعلاقته بوالديه. فقد يختلف إدراك المراهقين للممارسة الوالدية نفسها، ممّا يجعل
فهم المعنى الذي يمنحه المراهق لأسلوب معاملة والديه مدخلًا أكثر عمقًا لفهم الطريقة التي يبني بها نظرته
إلى ذاته ويقيّم بها قيمته الشخصية.

وانطلاقًا من الأهمية التي يحظى بها كلٌّ من أسلوب المعاملة الوالدية وتقدير الذات في تفسير النموّ النفسي
والاجتماعي للمراهق، فقد اهتمّ عددٌ كبير من الباحثين بدراسة العلاقة بين هذين المفهومين، انطلاقًا من
اختلاف الأساليب الوالدية وما قد تتركه من آثار متباينة في شخصية الأبناء. وقد تنوّعت هذه الدراسات من
حيث البيئات الثقافية التي أُجريت فيها، والمناهج والأدوات التي اعتمدتها، كما اختلفت نتائجها باختلاف
طبيعة العيّنات والمتغيّرات المدروسة، إلّا أنها اتّفقت على أنّ الأسرة تمثّل الإطار الأوّل الذي تتكوّن داخله
صورة المراهق عن ذاته، وأنّ أسلوب معاملة الوالدين يُعدّ من أهمّ العوامل المؤثّرة في هذا البناء النفسي.

وفي هذا الإطار، أجرت أيت مولود وإكردوشن بعلي (2018) دراسةً هدفت إلى التعرّف على العلاقة بين
أساليب المعاملة الوالدية وتقدير الذات لدى الأبناء المراهقين في ضوء الرتبة الميلادية. واعتمدت الدراسة
المنهج الوصفي الارتباطي، واستعملت أدواتٍ كمّية لقياس أساليب المعاملة الوالدية ومستوى تقدير الذات لدى
المراهقين. وأسفرت نتائجها عن وجود علاقة ذات دلالة بين أساليب المعاملة الوالدية وتقدير الذات، حيث
ارتبطت المعاملة القائمة على الدفء والتسامح والتقبّل بارتفاع مستوى تقدير الذات، في حين ارتبطت
الممارسات القائمة على العقاب والضبط الصارم والتمييز بين الأبناء بانخفاضه.

كما هدفت دراسة شكراوي (2025) إلى الكشف عن تأثير أساليب المعاملة الوالدية في تقدير الذات لدى
المراهقين، معتمدةً المنهج الوصفي، وشملت عيّنةً مكوّنةً من ثلاثين تلميذًا وتلميذة، واستعملت مقياس شيفر
\en{(Schaefer)} لأساليب المعاملة الوالدية، ومقياس كوبر سميث \en{(Coopersmith)} لتقدير الذات. وقد
توصّلت الدراسة إلى وجود علاقة واضحة بين المتغيّرين، حيث ارتبطت الأساليب الوالدية الإيجابية بارتفاع
مستوى تقدير الذات، بينما ارتبطت الأساليب السلبية وغير الفعّالة بانخفاضه، مؤكّدةً أنّ نوعية التفاعل
الأسري تؤدّي دورًا أساسيًا في تكوين الصورة التي يكوّنها المراهق عن ذاته.

أمّا الدراسة المرجعية التي أجراها لامبورن وزملاؤه \en{(Lamborn et al., 1991)} فقد هدفت إلى مقارنة
مستوى الكفاءة النفسية والاجتماعية لدى المراهقين في ضوء أنماط المعاملة الوالدية المختلفة، واعتمدت
المنهج الكمّي على عيّنة كبيرة تجاوزت أربعة آلاف مراهق، مستندةً إلى تصنيف بومريند الذي طوّره ماكوبي
ومارتن. وقد بيّنت نتائجها أنّ المراهقين المنتمين إلى أسرٍ تتبنّى النمط الديمقراطي يتمتّعون بمستويات أعلى
من الكفاءة النفسية والاستقلالية والتحصيل الدراسي وتقدير الذات، مقارنةً بالمراهقين الذين يعيشون في أسرٍ
يغلب عليها الأسلوب المتسلّط أو المُهمِل، وهو ما شكّل أحد أهمّ الأسس النظرية التي استندت إليها الدراسة
الحالية في تفسير أثر المعاملة الوالدية في بناء تقدير الذات.

وفي المقابل، قدّمت الدراسة الحديثة المُنجَزة في نيبال (2025) نتائج تستحقّ التوقّف عندها، إذ هدفت إلى
دراسة العلاقة بين أنماط المعاملة الوالدية وكلٍّ من الاكتئاب والقلق والضغط النفسي وتقدير الذات لدى
المراهقين، واعتمدت المنهج الكمّي، وشملت \en{583} مراهقًا ومراهقة، واستعملت استبيان أساليب السلطة
الوالدية ومقياس روزنبرغ \en{(Rosenberg)} لتقدير الذات، إضافةً إلى مقياس \en{DASS-21} لقياس
الاكتئاب والقلق والضغط النفسي. وأظهرت نتائجها أنّ العلاقة بين الأسلوب المتسلّط وتقدير الذات لم تكن
بالسلبية نفسها التي توصّلت إليها معظم الدراسات السابقة، بل ظهر ارتباط إيجابي ضعيف، وهي نتيجة
وصفها الباحثون بأنها غير متوقّعة، وأرجعوها إلى الخصائص الثقافية والاجتماعية للعيّنة، وإلى اختلاف
طريقة إدراك المراهقين للممارسات الوالدية داخل المجتمع النيبالي.

وبالنظر إلى هذه الدراسات السابقة، يتبيّن وجود اتّجاه عامّ يكاد يكون متّفقًا عليه، إذ انتهت ثلاث منها --
على الرغم من اختلاف البيئات التي أُجريت فيها والعيّنات التي اعتمدتها -- إلى أنّ إدراك المراهق لأسلوب
معاملة والديه المتسلّط يرتبط بانخفاض مستوى تقديره لذاته. ومع ذلك، لا يبدو هذا الاتّجاه مطلقًا، فقد جاءت
الدراسة النيبالية بنتيجة مختلفة تلفت الانتباه إلى أنّ فهم العلاقة بين أسلوب المعاملة الوالدية وتقدير الذات لا
يمكن أن ينفصل عن السياق الذي يعيش فيه المراهق، ولا عن الطريقة التي يدرك بها خبراته الأسرية
ويمنحها معنى.

وفي ضوء ذلك، لا تتمثّل الفجوة البحثية في غياب الدراسات التي تناولت العلاقة بين أسلوب المعاملة
الوالدية وتقدير الذات، وإنما في غياب دراسات كيفية -- وعلى وجه الخصوص في البيئة المغربية -- تستكشف
هذه العلاقة من منظور المراهق نفسه، وتكشف كيف يدرك أسلوب معاملة والديه المتسلّط، وكيف يفسّر أثر
هذا الإدراك في بناء تقديره لذاته عبر أبعاده المختلفة. ومن هنا جاءت الدراسة الحالية لتوظّف المقاربة
الكيفية، باعتبارها الأنسب لفهم الخبرة الذاتية للمراهق، والكشف عن المعاني التي يبنيها انطلاقًا من تجاربه
الأسرية، بعيدًا عن الاقتصار على تفسير العلاقة من خلال المؤشّرات الإحصائية وحدها.

\section{أهمية البحث}
تنبع أهمية هذا البحث من كونه يسعى إلى تقديم فهمٍ أعمق للعلاقة بين إدراك أسلوب المعاملة الوالدية
المتسلّط وتقدير الذات لدى المراهقين، من خلال الاقتراب من الخبرة كما يعيشها المراهق نفسه، لا كما
تعكسها المقاييس النفسية وحدها. كما يأمل أن يسهم في إثراء المكتبة النفسية العربية والمغربية بدراسة كيفية
في موضوعٍ قلّ تناوله بهذه المقاربة، وأن يوفّر للمهتمّين والمربّين قراءةً أعمق لكيفية إدراك المراهقين
لممارسات والديهم، بما قد يفيد في توجيه البرامج الإرشادية الأسرية. ويفتح البحث المجال أمام دراسات
مستقبلية تتناول متغيّرات أخرى مرتبطة بالمناخ الأسري، والقدوة الوالدية، والصراعات الأسرية، والفروق في
إدراك الإخوة داخل الأسرة الواحدة.
